% \documentclass[12pt]{article}

%%v helpful MIT exam schema, see http://www-math.mit.edu/~psh/exam/examdoc.pdf
 \documentclass[8pt]{exam} 
% \qformat{\textbf{Question \thequestion}\quad \thepoints\hfill } % Large depth to make space}
 \renewcommand{\thesubpart}{(\roman{subpart})}
 \renewcommand{\subpartlabel}{\thesubpart}
 \footer{}{}{Page \thepage\ de \numpages}
\addpoints
% \pointsinmargin
% \boxedpoints
\renewcommand{\questionshook}{\setlength{\itemsep}{0.2cm}}
\renewcommand{\partshook}{
	\setlength{\itemsep}{0.1cm}
}

\pagestyle{empty} %for handouts
\addtolength{\jot}{2em} % this makes all formulae a bit more spaced vertically for handouts
\setlength{\parindent}{0pt} %we're not writing a novel
%\renewcommand{\arraystretch}{2} %makes tables taller, so easier to write in

\usepackage[fleqn]{amsmath}
\usepackage{amssymb} %standard classy maths symbols
\usepackage[a4paper, margin=1in]{geometry} %makes it easy to specify where to put stuff on the page

\usepackage{siunitx} %v handy way of getting SI units to typeset correctly...
% \sisetup{locale = FR} %... and now it will use a comma for decimals

\usepackage[utf8]{inputenc}

% \usepackage[french]{babel} % frenchify everything (e.g. a Table becomes a Tableau)
% \usepackage{lmodern} % font with nicer French characters than standard
% \usepackage{textcomp} % and more help for special characters

\usepackage{tcolorbox} % basic but good package for boxes around text
\tcbset{colback=black!5!white}
% \usepackage{color} % colours
%\newcommand{\fillin}{\color{black!1!white}} %makes the text disappear
%\definecolor{fillin}{rgb} {1.00,1.00,1.00}
%\renewcommand{\fillin}{\color{blue}} \definecolor{fillin}{rgb} {0.00,0.00,1.00} %makes the filling in text appear (in blue)

\usepackage{enumitem} % can change the labels on lists, items, enumerates
\usepackage{setspace}

\usepackage{tikz, pgfplots} % interval diagrams, sets, etc.
\usetikzlibrary{calc,trees,positioning,arrows,fit,shapes}

\newcommand\TikCircle[1][2.5]{\tikz[baseline=-#1]{\draw[thick](0,0)circle(3mm)[radius=#1mm];}}

\begin{document}

\textbf{Name:} \dotfill

\

\begin{questions}

\question[3] Complete the statement of the theorem describing the conditions on rational roots:
\vspace{-5mm}
\fullwidth{
\begin{center}
\begin{minipage}{0.8\textwidth}
\setstretch{2.0}

Let $P(x) = a_n x^n + a_{n-1} x^{n-1} + \ldots + a_1 x + a_0$ be a polynomial with a rational root. We assume the root is $\frac{p}{q}$ in reduced form, which means $p,q \in \mathbb{Z}$ and
\fillwithdottedlines{1.3cm}
Then we have the following conditions on $p$ and $q$:
\fillwithdottedlines{1.3cm}

\end{minipage}
\end{center}
}

\question[2] Put numbers in the circles to indicate the correct order of this proof:

\begin{itemize}[label=\TikCircle]
\item $ a_n p^n = q \Bigg(- a_{n-1} p^{n-1} - \ldots - a_1 pq^{n-2} - a_0 q^{n-1} \Bigg)$
\item $ a_n \bigg(\dfrac{p}{q}\bigg)^n + a_{n-1} \bigg(\dfrac{p}{q}\bigg)^{n-1} + \ldots + a_1 \bigg(\dfrac{p}{q}\bigg) + a_0 = 0$
\item $P\bigg(\dfrac{p}{q}\bigg) = 0$
\item $ a_n p^n + a_{n-1} p^{n-1}q + \ldots + a_1 pq^{n-1}+ a_0 q^n = 0$
\item $ a_n \dfrac{p^n}{q^n} + a_{n-1} \dfrac{p^{n-1}}{q^{n-1}} + \ldots + a_1 \dfrac{p}{q} + a_0 = 0$
\end{itemize}

\question[2] Complete the proof to show the condition on $q$
\fillwithdottedlines{3cm}

\question[3] List the possible \textbf{positive} roots of $2x^{17} + x^3 - 49$ (but no need to check them).
\fillwithdottedlines{2cm}


\vfill

\end{questions}

\begin{tcolorbox}

\textbf{Name of examiner:}

\begin{center}
\gradetable[h][questions]
\end{center}

\end{tcolorbox}


\end{document}